% preface.tex — Book preface
\chapter*{Preface}
\addcontentsline{toc}{chapter}{Preface}

This book is a guided tour through the physics, mathematics, and numerical
methods needed to understand---and contribute to---the \texttt{event-horizon}
codebase, a relativistic ray tracer for black hole visualisation. Rather than
treating general relativity and numerical computation as separate disciplines,
we develop them together, showing how every equation maps to running code.

\paragraph{Prerequisites.}
We assume the reader has completed a standard course in special relativity and
is comfortable with linear algebra, multivariable calculus, and basic
differential equations. No prior exposure to general relativity or numerical
methods is required.

\paragraph{Structure.}
The book is organised in seven parts:
\begin{enumerate}
  \item \textbf{Foundations} --- special relativity, differential geometry, and
        the Einstein equations.
  \item \textbf{Black Hole Spacetimes} --- Schwarzschild, Kerr, and geodesic
        motion.
  \item \textbf{Numerical Methods} --- finite differences, ODE integration,
        automatic differentiation, and spectral methods.
  \item \textbf{Numerical Relativity} --- the ADM formalism and the BSSN
        evolution system.
  \item \textbf{GRMHD} --- general-relativistic magnetohydrodynamics and
        conservative schemes.
  \item \textbf{Visualisation} --- ray tracing, radiative transfer, and
        rendering.
  \item \textbf{Advanced Topics} --- extensions, optimisations, and open
        problems.
\end{enumerate}

\paragraph{Notation.}
We use the $({-}{+}{+}{+})$ metric signature throughout. Greek indices
$\mu,\nu,\ldots$ run over spacetime dimensions $0$--$3$; Latin indices
$i,j,\ldots$ run over spatial dimensions $1$--$3$. We work in geometrised
units where $G = c = 1$ unless otherwise stated. Sections marked with a
$\star$ are optional or advanced and may be skipped on a first reading.
